\documentclass[a4paper,11pt]{article}

\usepackage{latexsym}
\usepackage{xcolor}
\usepackage{float}
\usepackage{ragged2e}
\usepackage[empty]{fullpage}
\usepackage{wrapfig}
\usepackage{lipsum}
\usepackage{tabularx}
\usepackage{titlesec}
\usepackage{geometry}
\usepackage{marvosym}
\usepackage{verbatim}
\usepackage{enumitem}
\usepackage[hidelinks]{hyperref}
\usepackage{fancyhdr}
\usepackage{multicol}
\usepackage{graphicx}
\usepackage{cfr-lm}
\usepackage[T1]{fontenc}
\usepackage{fontawesome}
\usepackage[most]{tcolorbox}

\geometry{left=0.85cm, top=0.8cm, right=0.85cm, bottom=0.2cm}
\titleformat{\section}{\vspace{-4pt}\scshape\raggedright\large}{}{0em}{}[\color{black}\titlerule \vspace{-7pt}]

\newcommand{\resumePOR}[3]{%
\vspace{0.5mm}\item[]
    \begin{tabular*}{\textwidth}[t]{l@{\extracolsep{\fill}}r}
    \hspace{-3mm}{#1}:\hspace{1mm} & \hspace*{0pt}\hfill{\footnotesize{ #3}} \vspace{-0.5mm}\\ \hspace{-2.9mm}#2
    \end{tabular*}
    \vspace{0mm}
}
\newcommand{\resumeExp}[4]{%
\vspace{2mm}\item[]
    \begin{tabular*}{\textwidth}[t]{l@{\extracolsep{\fill}}r}
        \hspace{-4.4mm} \small\textbf{#1} & {\footnotesize{#3}}\vspace{-1.2mm}\\
        \hspace{-4.3mm} \footnotesize{\text{#2}} & \footnotesize{#4}
    \end{tabular*}
    \vspace{-3.2mm}
}
\newcommand{\resumeProject}[4]{%
\vspace{2mm}\item[]
    \begin{tabular*}{\textwidth}[t]{l@{\extracolsep{\fill}}r}
        \hspace{-4.4mm} \small\textbf{#1} & {\footnotesize{#3}}\vspace{-1mm}\\
        \hspace{-4.3mm} \footnotesize{\text{#2}} & \footnotesize{#4}
    \end{tabular*}
    \vspace{-3.2mm}
}
\newcommand{\resumeEdu}[4]{%
\vspace{0mm}\item[]
    \begin{tabular*}{\textwidth}[t]{l@{\extracolsep{\fill}}r}
        \hspace{-4.3mm} \small\textbf{#1} & \footnotesize{#3}\vspace{-1mm} \\
        \hspace{-4.3mm} \footnotesize{#2} & \footnotesize{#4}
    \end{tabular*}
    \vspace{-3.2mm}
}
\newcommand{\resumeAchieve}[3]{%
\hspace{-3.1mm}\textbf{ #1} & {#2} & \hspace{3mm}\footnotesize{#3}\vspace{0mm}\\
}
\renewcommand{\labelitemi}{$\vcenter{\hbox{\tiny$\bullet$}}$}
\newcommand{\resumeSubHeadingListStart}{\begin{itemize}[leftmargin=*,labelsep=0mm,itemsep=-2.5mm]}
\newcommand{\resumeItemListStart}{\begin{justify}\begin{itemize}[leftmargin=3ex, rightmargin=2ex, noitemsep,labelsep=1.2mm,itemsep=0mm]\small}
\newcommand{\resumeSubHeadingListEnd}{\end{itemize}\vspace{-2mm}}
\newcommand{\resumeItemListEnd}{\end{itemize}\end{justify}\vspace{-1.5mm}}
\renewcommand{\arraystretch}{1}
\newcolumntype{L}{>{\raggedright\arraybackslash}X}
\newcolumntype{R}{>{\raggedleft\arraybackslash}X}
\newcolumntype{C}{>{\centering\arraybackslash}X}

\newcommand{\name}{Jair Santana}
\newcommand{\phone}{+1-857-799-0252}
\newcommand{\emaila}{jairrs@mit.edu}
\newcommand{\github}{rekomodo}
\newcommand{\linkedin}{jairrs}

% Content toggles. Uncomment a \def line to include that entry.

% Education
\def\showMITEducation{}

% Experience
\def\showSalesforceExperience{}
\def\showCapcomExperience{}
%\def\showFESTIMResearch{}

% Projects
\def\showRaftProject{}
\def\showSokobamProject{}
\def\showFPGAProject{}
%\def\showProcessorDesignProject{}
%\def\showRealTimeGraphicsProject{}
%\def\showFMRadiopcbProject{}
%\def\showOtherProjectsEntry{}

% Achievements
\def\showIberoamericanOlympiadAchievement{}
\def\showInternationalOlympiadAchievement{}
\def\showDominicanOlympiadAchievement{}

% Positions of Responsibility
\def\showSoftwareConstructionTA{}
\def\showDominicanOlympiadTrainer{}

\begin{document}
\fontfamily{cmr}\selectfont

\begin{center}
    \LARGE{\textbf{\name}}
\end{center}
\vspace{-4.5mm}
\begin{center}
    \small{\href{https://github.com/\github}{\faGithub \hspace{0.2mm} github.com/\github} | \href{https://www.linkedin.com/in/\linkedin/}{\faLinkedinSquare \hspace{0.2mm} linkedin.com/in/\linkedin} | \href{mailto:\emaila}{\faSend \hspace{0.2mm} \emaila} | \faPhone \hspace{0.2mm} \phone}
\end{center}
\vspace{-3mm}

% Comment or uncomment section inputs to choose sections.
\vspace{-2.5mm}
\section{Education}
\resumeSubHeadingListStart
\ifdefined\showMITEducation
\resumeEdu
{Massachusetts Institute of Technology, Cambridge MA}
{Bachelor of Science in Computer Science and Engineering}
{May 2027}
{GPA: 4.9/5.0}
\resumeSubHeadingListEnd
\fi
\vspace{-3.5mm}

\section{Relevant Coursework}
\vspace{0.2mm}
\small{\begin{tabular*}{\textwidth}[t]{p{\textwidth}}
\hspace{-3mm}\textbf{ Past Coursework: }{Advanced Algorithms, Operating Systems, FPGA Engineering, Theory of Computation}\\
\hspace{-3mm}\textbf{ Current Coursework: }{Distributed Systems, Computer Systems Security, Secure Hardware Design}
\end{tabular*}}
\vspace{-4mm}

\section{Experience}
\resumeSubHeadingListStart
\ifdefined\showSalesforceExperience
\resumeExp
{Salesforce - Field Service Mobile (SFS Mobile Comms \& Search)}
{Software Engineer Intern}
{Summer 2026}
{San Francisco, CA}
\resumeItemListStart
\item[$\bullet$] Shipped my first customer-facing fix within 7 days of joining and merged 19 pull requests---approximately 63\% of the average contributor's PR count over the same period on a team of senior engineers---after rapidly ramping into a cross-platform iOS and Android codebase with no prior mobile-development experience.
\item[$\bullet$] Resolved 33 work items, including 15 fixes delivered through merged code, and closed 6 accessibility defects within SLO across iOS and Android, supporting WCAG compliance and public-sector eligibility.
\item[$\bullet$] Migrated a legacy screen to SwiftUI and shared design-system components, completing the iOS implementation and advancing the corresponding Android implementation to code review.
\item[$\bullet$] Identified and filed 8 reproducible defects across iOS and Android during v264 regression testing, including an iOS-only crash caught before release.
\resumeItemListEnd
\fi

\ifdefined\showCapcomExperience
\resumeExp
{CAPCOM - Game Engine Team}
{Game Engine Development Intern}
{May 2025 - Aug 2025}
{Osaka, Japan}
\resumeItemListStart
\item[$\bullet$] Architected and implemented a modular C\# tooling framework integrated with proprietary C++ engine APIs, enabling live hot-reloading of developer-authored scripts without full engine recompilation.
\item[$\bullet$] Designed custom APIs, hot-reloading schemes, data-collection pipelines, and supporting engine systems to make asset data more flexible for developers and improve iteration speed.
\item[$\bullet$] Designed and implemented asset-analysis algorithms enabling real-time introspection of engine data, reducing implementation iteration time for game developers.
% \item[$\bullet$] Designed a WPF-based MVVM UI layer to support dynamic data-binding and real-time engine state visualization.
\resumeItemListEnd
\fi

\ifdefined\showFESTIMResearch
\resumeExp
{Plasma Science and Fusion Center - FESTIM}
{Verification and Validation Developer, Research Assistant}
{June 2024 - Aug 2024}
{Cambridge, MA}
\resumeItemListStart
\item[$\bullet$] Developed reproducible verification and validation benchmarks for \href{https://festim.readthedocs.io/}{FESTIM}, publishing automated test suites via Jupyter Book to ensure scientific reproducibility.
\item[$\bullet$] Refactored core components for improved stability and maintainability.
\item[$\bullet$] Reviewed scientific literature to identify new validation cases and integrated them into the Jupyter Book.
\resumeItemListEnd
\fi

\ifdefined\showSoftwareConstructionTA
\resumePOR{\textbf{Software Construction Undergraduate TA, MIT}}
{Guide students towards developing software development best-practices and develop course materials.}
{\raisebox{0.75pt}{Jan 2026 - May 2026}}
\fi
\vspace{0.5mm}
\ifdefined\showDominicanOlympiadTrainer
\resumePOR{\textbf{Team Trainer and Problem Setter, Dominican Informatics Olympiad (ODI)}}
{Prepare lectures for the national team on DSA and competitive programming techniques.}{\raisebox{0.75pt}{2021 - Present}}
\fi
\resumeSubHeadingListEnd
\vspace{-3.5mm}

\section{Projects}
\resumeSubHeadingListStart
\ifdefined\showRaftProject
\resumeProject
{Raft Consensus Algorithm}
{Distributed Systems, MIT 6.5840}
{Feb 2026 - May 2026}
{}
\resumeItemListStart
\item[$\bullet$] Implemented a fault-tolerant consensus algorithm in Go, featuring leader election, log replication, and safety guarantees under extreme network partitions and server failures.
\item[$\bullet$] Engineered a robust state machine supporting persistent storage and snapshotting to manage memory constraints and accelerate recovery during node reboots.
\item[$\bullet$] Optimized concurrency using Go’s synchronization primitives (mutexes, condition variables) to handle high-throughput RPC requests while avoiding race conditions and deadlocks.
\resumeItemListEnd
\fi

\ifdefined\showSokobamProject
\resumeProject
{Sokobam {\href{https://github.com/rekomodo/sokobam}{\footnotesize\faExternalLink \hspace{0.1mm}}}}
{Real-Time Multiplayer Game}
{March 2026 - Present}
{}
\resumeItemListStart
\item[$\bullet$] Developed a live, web-based multiplayer puzzle game using React and SpacetimeDB, allowing multiple players to interact in a shared persistent world with low-latency state updates.
\item[$\bullet$] Implemented client-side rendering and real-time synchronization of player game states.
\resumeItemListEnd
\fi

\ifdefined\showFPGAProject
\resumeProject
{FPGA Design}
{FPGA Vector Operations Module, MIT 6.205}
{Oct 2025 - Dec 2025}
{}
\resumeItemListStart
\item[$\bullet$] Designed and implemented a fully pipelined hardware vector normalizer in VHDL for AMD's Urbana development board, specialized for high-throughput repeated vector operations for use in graphics rendering.
\item[$\bullet$] Optimized for precision and speed by using numerical approximations where necessary, without compromising result quality, and reducing the module's resource usage to fit within project constraints.
\item[$\bullet$] Built a cocotb-based verification suite in Python, achieving full functional coverage across randomized vector test cases.
\resumeItemListEnd
\fi
\resumeSubHeadingListEnd
\vspace{-3.5mm}

\ifdefined\showProcessorDesignProject
\resumeSubHeadingListStart
\resumeProject
{Processor Design}
{Specialized RISC-V Processor, MIT}
{Apr 2020 - Jul 2020}
{}
\resumeItemListStart
\item[$\bullet$] Implemented a RISC-V processor in MIT's Minispec HDL, specialized for high-throughput repeated vector operations in neural network training workloads.
\item[$\bullet$] Designed architectural optimizations including pipelining, cache modifications, and custom instructions.
\resumeItemListEnd
\resumeSubHeadingListEnd
\fi

\ifdefined\showRealTimeGraphicsProject
\resumeSubHeadingListStart
\resumeProject
{6.440 Final Project}
{Real-Time Graphics Project, MIT}
{Sep 2025 - Dec 2025}
{\href{https://github.com/rekomodo/6440_final}{github.com/rekomodo/6440_final}}
\resumeItemListStart
\item[$\bullet$] Implemented a real-time interactive graphics application in Godot, focusing on custom rendering behavior and scene-level control.
\item[$\bullet$] Developed and integrated custom shaders (GLSL/Godot shader language) to control lighting, materials, and visual effects.
\item[$\bullet$] Worked directly with rendering pipeline concepts including transformations, coordinate spaces, and per-frame update logic.
\item[$\bullet$] Tuned performance by analyzing draw calls, shader cost, and scene complexity to maintain real-time responsiveness.
\resumeItemListEnd
\resumeSubHeadingListEnd
\fi

\ifdefined\showFMRadiopcbProject
\resumeSubHeadingListStart
\resumeProject
{PCB Design}
{FM Radio, MIT}
{Jan 2026}
{}
\resumeItemListStart
\item[$\bullet$] Laid-out and routed the design for an FM Radio via KiCad.
\item[$\bullet$] Designed audio amplifier and low-pass filter components and verified them in LTSpice.
\item[$\bullet$] Soldered components onto printed board and debugged for problems.
\resumeItemListEnd
\resumeSubHeadingListEnd
\fi

\ifdefined\showOtherProjectsEntry
\resumeSubHeadingListStart
\resumeProject{Other Projects}{}{}{}
\resumeItemListStart
\item[$\bullet$] \textbf{Sokobam:} Live multiplayer web game implemented with React and SpacetimeDB.
\item[$\bullet$] \textbf{Raft Consensus Algorithm:} Implementation of Raft in Go for MIT's Distributed Systems class.
\resumeItemListEnd
\resumeSubHeadingListEnd
\fi

\section{Skills}
\vspace{0.2mm}
\small{\begin{tabular*}{\textwidth}[t]{p{\textwidth}}
\hspace{-2mm}\textbf{Programming Languages:} C, C++, C\#, TypeScript, Go, Python, Rust, RISC-V Assembly \\
\hspace{-2mm}\textbf{Systems \& Hardware:} SystemVerilog, cocotb, FPGA Design, KiCad, LTSpice \\
\hspace{-2mm}\textbf{Tools:} Git, Linux, Bash, Docker, GLFW, GLSL, WPF, Cursor \\
\hspace{-2mm}\textbf{Languages:} Spanish (Native), English (Fluent), Japanese (Intermediate)
\end{tabular*}}

\vspace{-2mm}
\section{Achievements}
\vspace{0.2mm}
\small{\begin{tabular*}{\textwidth}[t]{p{0.4\textwidth} p{0.6\textwidth}@{\extracolsep{\fill}}r}
\ifdefined\showIberoamericanOlympiadAchievement
\resumeAchieve{Iberoamerican Informatics Olympiad (OII)}{Silver Medalist}{2021}
\fi
\ifdefined\showInternationalOlympiadAchievement
\resumeAchieve{International Informatics Olympiad (IOI)}{Three-Time Participant}{2021-2023}
\fi
\ifdefined\showDominicanOlympiadAchievement
\resumeAchieve{Dominican Informatics Olympiad (ODI)}{First Perfect Score; Three-Time Gold Medalist}{2020-2023}
\fi
\end{tabular*}}
\vspace{-2mm}

%\section{Positions of Responsibility}
\vspace{-0.4mm}
\resumeSubHeadingListStart
\ifdefined\showSoftwareConstructionTA
\resumePOR{\textbf{Software Construction Undergraduate TA, MIT}}
{Guide students towards developing software development best-practices and develop course materials.}
{\raisebox{0.75pt}{Jan 2026 - Present}}
\fi
\vspace{0.5mm}
\ifdefined\showDominicanOlympiadTrainer
\resumePOR{\textbf{Team Trainer and Problem Setter, Dominican Informatics Olympiad (ODI)}}
{Prepare lectures for the national team on DSA and competitive programming techniques.}{\raisebox{0.75pt}{2021 - Present}}
\fi
\vspace{-1mm}
\resumeSubHeadingListEnd
\hspace*{-2mm}\rule{1.030\textwidth}{0.1mm}
\vspace{0mm}


\end{document}
